% Bewertungspaket zum Gesamttest «Energie» — Quelle des PDFs (python3 scripts/build-lp-pdf.py energie).
% Ganze Punkte je Teilschritt; Werte und Folgefehler-Fälle mit python3 nachgerechnet (06.10.2026).
% Fassung 3 (06.10.2026): G2 bis G4 und G6d neu, gleichartige Fehler gleich bewertet, «ein Fehler,
% ein Abzug» im Auftrag an die KI, Tabellenzeilen mit mehr Abstand (Brüche berühren sich nicht).
% Aufbau, Auftrag an die KI und Punktarten (E, W, B, A) wie im Leitprogramm Dynamik.
\documentclass[11pt]{article}
\usepackage{../lp-druck}
\renewcommand{\lpname}{Energie}
\renewcommand{\lpteilgebiet}{4.3}
\renewcommand{\lpfuss}{Bewertungspaket}
\newcolumntype{P}{>{\centering\arraybackslash}p{10mm}}
\newenvironment{raster}{\par\smallskip\noindent\renewcommand{\arraystretch}{1.5}\tabularx{\linewidth}{@{}>{\raggedright\arraybackslash}p{38mm}P >{\raggedright\arraybackslash}X@{}}\toprule
  \small\textsc{Teilschritt} & \small\textsc{P} & \small\textsc{Musterlösung}\\\midrule}{\bottomrule\endtabularx\par}
\newcommand{\fehler}[1]{\par\smallskip{\small\textcolor{lprot}{\bfseries Typische Fehler:} #1}}
\newcommand{\E}{\,{\footnotesize\bfseries\color{lpgruen}(E)}}
\newcommand{\B}{\,{\footnotesize\bfseries(B)}}
\newcommand{\W}{\,{\footnotesize\bfseries(W)}}
\definecolor{lplila}{HTML}{5B2D8E}
\newcommand{\A}{\,{\footnotesize\bfseries\color{lpgruen}(A)}}
\begin{document}
\lpkopf{Bewertungspaket zum Gesamttest}{}

\begin{lpachtung}{Erst lösen, dann öffnen}
Dieses Paket enthält die Musterlösung. Wer es vor dem Lösen liest, misst am Ende nur, wie gut das Abschreiben gelingt.
\end{lpachtung}

\section*{1 · So gehst du vor}
\begin{enumerate}[leftmargin=*,itemsep=2pt]
\item \textbf{Gesamttest lösen} — vollständig, mit Rechenweg, ohne dieses Paket.
\item \textbf{Fotografieren oder scannen}, gut lesbar, eine Seite pro Bild. \textbf{Kein Name} auf den Bildern.
  Handyfotos tragen oft unsichtbar Standort, Zeit und Gerät mit (Metadaten). Lade darum lieber einen Scan oder ein
  Bildschirmfoto des Fotos hoch, oder entferne vorher die Standortangaben.
\item \textbf{Bewerten lassen.} Ob du dafür eine KI nutzt und welche, entscheidest du selbst und in eigener Verantwortung —
  je nachdem, was dir aufgrund deines Alters und deiner persönlichen Situation erlaubt ist (Altersgrenzen und
  Nutzungsbedingungen des Anbieters, allenfalls Einverständnis deiner Eltern). Lade nur deine Lösung und dieses PDF hoch,
  keine persönlichen Daten, und füge den Auftrag aus Abschnitt~2 ein. Ohne KI kannst du dich mit Abschnitt~3 selbst bewerten.
\item \textbf{Rückmeldung prüfen.} Stimmt, was die KI aus deiner Handschrift gelesen hat? Eine KI kann sich irren —
  im Zweifel gilt das Raster in Abschnitt~3.
\item \textbf{Gezielt wiederholen} nach Abschnitt~4.
\end{enumerate}

\needspace{14\baselineskip}
\section*{2 · Auftrag an die KI}
\begin{tcolorbox}[colback=lplinie!25,colframe=lplinie,boxrule=0.5pt,arc=1mm,fontupper=\small]
Du bist Korrektorin oder Korrektor für Physik an einer Schweizer Berufsmaturitätsschule. Im Anhang findest du (1)~das Bewertungspaket zum Gesamttest «Energie» mit Musterlösung und Punkteraster und (2)~Fotos meiner handschriftlichen Lösung.\par\medskip
Geh so vor:\par
1. Schreib zu jeder Aufgabe G1 bis G6 zuerst ab, was du in meiner Lösung liest — Rechenweg und Ergebnis. Was du nicht sicher lesen kannst, markierst du mit [unsicher]. Nicht raten.\par
2. Vergib die Punkte genau nach dem Raster im Bewertungspaket (Abschnitt 3), ganze Punkte je Zeile. Ziehe einen Fehler nur einmal ab: Rechne ich mit einem falschen Zwischenergebnis richtig weiter, gibt es die Folgepunkte — auch wenn das Folgeergebnis unplausibel ist; eine fehlende Plausibilitätsprobe kostet keinen zusätzlichen Punkt. Derselbe Fehler in mehreren Teilaufgaben (z.\,B. Taschenrechner im Bogenmass, sin/cos vertauscht, dieselbe fehlende Umrechnung) kostet nur beim ersten Vorkommen einen Punkt; die anderen Teilaufgaben werden folgerichtig bewertet. Gleichartige Fehler werden gleich bewertet. Die Zeilen «Typische Fehler» sagen, welche Rasterzeile bei diesem Fehler 0 Punkte gibt — sie sind kein zusätzlicher Abzug.\par
3. Andere richtige Lösungswege zählen voll. Gleichwertige Schreibweisen zählen voll: Dezimalpunkt oder Dezimalkomma, andere Einheit oder Vorsilbe ($0.12\;\text{m} = 12\;\text{cm}$; $\text{m/s}$ statt $\text{km/h}$, wenn richtig umgerechnet), Zehnerpotenz oder ausgeschriebene Zahl, sinnvoll gerundete Werte (drei signifikante Stellen). Zeilen mit (E) sind Ergebnispunkte: Sie verlangen das Ergebnis \emph{mit richtiger Einheit}, auch ohne Weg; ohne oder mit falscher Einheit gibt es diesen Punkt nicht. Zeilen mit (W) gibt es nur mit nachvollziehbarem Weg (Formel, dann Werte mit Einheiten); ein Zahlenwert darin braucht ebenfalls die Einheit. Zeilen mit (B) sind Begründungen: Es zählt die fachliche Aussage in eigenen Worten, eine Formel ist nicht nötig; andere richtige Formulierungen zählen voll. Zeilen mit (A) sind Ablesepunkte: der richtig aus dem Diagramm abgelesene Wert mit Einheit. Werte aus sinnvoll gerundeten Zwischenergebnissen zählen, wenn sie höchstens rund 2\,\% vom Raster abweichen.\par
4. Begründe jeden Abzug in einem Satz und nenne die Fehlerart (zum Beispiel «Wurzel vergessen»). Schreib die Musterlösung nicht ab — zeig mir, wo mein Fehler liegt.\par
5. Gib zum Schluss eine Tabelle «Aufgabe | Punkte | Begründung», die Gesamtpunktzahl (von 25) und — nach der Selbsteinschätzung in Abschnitt 4 — welches Kapitel des Leitprogramms ich wiederholen soll. Keine Note.\par
6. Fehlt ein Foto oder ist eine Aufgabe unlesbar, sag es, statt sie zu bewerten.
\end{tcolorbox}

\needspace{16\baselineskip}
\section*{3 · Musterlösung und Punkteraster}
\textbf{Allgemein:} Ganze Punkte je Zeile. Ein Fehler wird nur einmal abgezogen: Wer mit einem falschen Zwischenergebnis
richtig weiterrechnet, bekommt die folgenden Zeilen (Folgefehler) — auch wenn das Folgeergebnis unplausibel ist; eine fehlende
Plausibilitätsprobe kostet keinen zusätzlichen Punkt. Derselbe Fehler in mehreren Teilaufgaben (z.\,B. Taschenrechner im Bogenmass,
sin/cos vertauscht, dieselbe fehlende Umrechnung) kostet nur beim ersten Vorkommen einen Punkt; die anderen Teilaufgaben werden folgerichtig
bewertet. Gleichartige Fehler werden gleich bewertet: Ein falscher Ansatz kostet den Ansatzpunkt und den Ergebnispunkt, der direkt daraus folgt;
alle weiteren Zeilen werden folgerichtig bewertet. Gleichwertige Wege und Schreibweisen zählen voll.
In der Spalte P: \textbf{(E)}~= Ergebnispunkt — es zählt das Ergebnis \emph{mit richtiger Einheit}, auch ohne Weg; ohne oder mit
falscher Einheit gibt es ihn nicht. \textbf{(W)}~= Wegpunkt — nur mit nachvollziehbarem Weg (Formel, dann Werte mit Einheiten);
steht in einer Weg-Zeile ein Zahlenwert, braucht auch er die Einheit. \textbf{(B)}~= Begründungspunkt — es zählt die fachliche
Aussage in eigenen Worten, eine Formel ist nicht nötig. \textbf{(A)}~= Ablesepunkt — der richtig abgelesene Wert mit Einheit.
Folgewerte aus gerundeten Zwischenergebnissen zählen, wenn sie höchstens rund 2\,\% abweichen. «Typische Fehler» nennt die Zeilen mit 0 Punkten und die Punktzahl, die bleibt.

\begin{lpaufgabe}{G1}{Begriffe}{4 P}
\begin{raster}
a) Energie & 1\B & die Fähigkeit, Arbeit zu verrichten (gleichwertig: Wärme abzugeben, Licht auszusenden), \emph{und} drei Formen, z.\,B. Lage-, Bewegungs-, Spannenergie, Wärme, chemische, elektrische, Strahlungs-, Kernenergie. Der Punkt verlangt beides; fehlt die Erklärung oder sind es weniger als drei Formen, gibt es ihn nicht\\
b) Energie und Leistung & 1\B & Leistung ist Energie (Arbeit) pro Zeit, \(P = \tfrac{W}{t}\); Beispiel, in dem dieselbe Energie schneller oder langsamer umgesetzt wird (z.\,B. Treppe rennen oder gehen; ein Wasserkocher mit 2 kW und einer mit 1 kW bringen dasselbe Wasser zum Sieden, brauchen aber verschieden lange)\\
c) Energieeffizienz & 1\B & dieselbe Nutzung mit möglichst wenig zugeführter Energie, also hoher Wirkungsgrad; Beispiel z.\,B. LED statt Glühlampe, Elektro- statt Benzinmotor, Wärmepumpe\\
d) Skifahrer & 1\B & Die Lageenergie nimmt ab, die Bewegungsenergie bleibt gleich (konstantes Tempo), die Wärme (Ski, Schnee, Luft) nimmt zu; die abgegebene Lageenergie wird durch Reibung und Luftwiderstand ganz zu Wärme (alle drei mit Begründung nötig)\\
\end{raster}
\fehler{a) nur Formen ohne Erklärung oder «Energie ist Kraft»: a) 0 \(\to\) 3 von 4\,P. b) «Leistung ist die Energie»: b) 0 \(\to\) 3 von 4\,P. d) «Die Bewegungsenergie nimmt zu» oder «Die Energie geht verloren»: d) 0 \(\to\) 3 von 4\,P.}
\end{lpaufgabe}

\begin{lpaufgabe}{G2}{Paket}{3 P}
\begin{raster}
a) Ablesen & 1\A & am Anfang \(F_s = 120\;\text{N}\), am Ende (bei \(8\;\text{m}\)) \(40\;\text{N}\)\\
a) Arbeit & 1\E & Fläche des Trapezes, z.\,B. Rechteck plus Dreieck: \(W = 40\;\text{N} \cdot 8\;\text{m} + \tfrac12 \cdot 80\;\text{N} \cdot 8\;\text{m} = 320\;\text{J} + 320\;\text{J} = 640\;\text{J}\) (gleichwertig: mittlere Kraft \(80\;\text{N}\) mal \(8\;\text{m}\))\\
b) Winkel & 1\E & \(W = F \cdot s \cdot \cos\alpha\), also \(\cos\alpha = \dfrac{W}{F \cdot s} = \dfrac{480\;\text{J}}{50\;\text{N} \cdot 12\;\text{m}} = 0.8\), \(\alpha \approx 36.9^\circ\)\\
\end{raster}
\fehler{a) Rechteck \(120\;\text{N} \cdot 8\;\text{m} = 960\;\text{J}\) oder Dreieck bis null \(\tfrac12 \cdot 120\;\text{N} \cdot 8\;\text{m} = 480\;\text{J}\): a) Arbeit 0 \(\to\) 2 von 3\,P. b) Sinus statt Kosinus (\(\alpha \approx 53.1^\circ\)) oder Taschenrechner im Bogenmass (\(0.644\)): b) 0 \(\to\) 2 von 3\,P.}
\end{lpaufgabe}

\begin{lpaufgabe}{G3}{Rollen und Bremsen: \(80\;\text{kg}\), \(6\;\text{m/s}\), \(8\;\text{m}\) hinunter, \(400\;\text{N}\)}{5 P}
\begin{raster}
Ansatz & 1\W & Energieerhaltung mit beiden Energien oben: \(m \cdot g \cdot h + \tfrac12 \cdot m \cdot v_0^2 = \tfrac12 \cdot m \cdot v^2\), Masse kürzen (gleichwertig: direkt \(v = \sqrt{v_0^2 + 2 \cdot g \cdot h}\) mit Werten)\\
a) Tempo unten & 1\E & \(v = \sqrt{(6\;\text{m/s})^2 + 2 \cdot 9.81\;\text{m/s}^2 \cdot 8\;\text{m}} \approx 13.9\;\text{m/s}\)\\
b) Bremsweg & 1\E & \(F_B \cdot s = \tfrac12 \cdot m \cdot v^2\), \(s = \dfrac{\tfrac12 \cdot 80\;\text{kg} \cdot (13.9\;\text{m/s})^2}{400\;\text{N}} \approx \dfrac{7718\;\text{J}}{400\;\text{N}} \approx 19.3\;\text{m}\) (gleichwertig: \(s = \dfrac{m \cdot g \cdot h + \tfrac12 \cdot m \cdot v_0^2}{F_B}\))\\
c) doppelt so schnell & 1\B & viermal so lang, rund \(77\;\text{m}\): Die Bewegungsenergie wächst mit \(v^2\), und dieselbe Bremskraft braucht für die vierfache Energie den vierfachen Weg\\
d) Form der Strasse & 1\B & nein: Ohne Reibung zählen in der Energiebilanz nur die Höhen (Start- und Endhöhe), nicht die Form oder Steilheit des Weges\\
\end{raster}
\fehler{Geschwindigkeiten addiert, \(v = 6\;\text{m/s} + \sqrt{2 \cdot g \cdot h} \approx 18.5\;\text{m/s}\), oder Anfangstempo vergessen, \(v = \sqrt{2 \cdot g \cdot h} \approx 12.5\;\text{m/s}\): Beides ist ein falscher Ansatz, gleich bewertet — Ansatz und a) 0, b) folgerichtig (\(34.3\;\text{m}\) bzw. \(15.7\;\text{m}\)) zählt \(\to\) 3 von 5\,P.
b) Faktor \(\tfrac12\) vergessen (\(38.6\;\text{m}\)): b) 0 \(\to\) 4 von 5\,P. c) «doppelt so lang»: c) 0 \(\to\) 4 von 5\,P.}
\end{lpaufgabe}

\begin{lpaufgabe}{G4}{Skilift: \(70\;\text{kg}\), \(500\;\text{m}\) Spur, \(120\;\text{m}\) hinauf, Seil \(200\;\text{N}\)}{4 P}
\begin{raster}
a) Seilarbeit & 1\E & \(W = F \cdot s = 200\;\text{N} \cdot 500\;\text{m} = 100\,000\;\text{J} = 100\;\text{kJ}\) (Seil und Weg in derselben Richtung, \(\cos 0^\circ = 1\))\\
b) Ansatz & 1\W & konstantes Tempo, die Bewegungsenergie bleibt: \(W = m \cdot g \cdot h + F_R \cdot s\), also \(F_R = \dfrac{W - m \cdot g \cdot h}{s}\) (gleichwertig: Kräfte längs der Spur, \(F_R = F - m \cdot g \cdot \tfrac{h}{s}\))\\
b) Reibung & 1\E & \(F_R = \dfrac{100\,000\;\text{J} - 82\,404\;\text{J}}{500\;\text{m}} \approx 35.2\;\text{N}\)\\
c) wohin & 1\B & beide Teile mit Betrag: rund \(82.4\;\text{kJ}\) in Lageenergie, der Rest, rund \(17.6\;\text{kJ}\), durch Reibung in Wärme (Ski und Schnee); die Bewegungsenergie ändert sich nicht. Folgerichtige Beträge aus einem falschen \(F_R\) zählen\\
\end{raster}
\fehler{Lageenergie vergessen (\(F_R = 200\;\text{N}\)), Lageenergie addiert statt abgezogen (\(F_R \approx 365\;\text{N}\)) oder Hubarbeit mit der Spurlänge statt der Höhe (negatives \(F_R\)): jedes Mal ein falscher Ansatz, gleich bewertet — b) Ansatz und b) Reibung 0 \(\to\) 2 von 4\,P.}
\end{lpaufgabe}

\begin{lpaufgabe}{G5}{Speicherkraftwerk: \(2000\;\text{kg}\) je Sekunde, \(300\;\text{m}\), \(\eta = 0.85\)}{4 P}
\begin{raster}
Ansatz & 1\W & \(P = \dfrac{W}{t} = \dfrac{m \cdot g \cdot h}{t}\), dann \(P_\text{el} = \eta \cdot P\)\\
Wasser & 1\E & \(P = \dfrac{2000\;\text{kg} \cdot 9.81\;\text{m/s}^2 \cdot 300\;\text{m}}{1\;\text{s}} \approx 5.89 \cdot 10^{6}\;\text{W} = 5.89\;\text{MW}\)\\
elektrisch & 1\E & \(P_\text{el} = 0.85 \cdot 5.89\;\text{MW} \approx 5.00\;\text{MW}\)\\
Energie in \(2\;\text{h}\) & 1\E & \(E = P_\text{el} \cdot t = 5.00\;\text{MW} \cdot 2\;\text{h} \approx 10.0\;\text{MWh}\); gleichwertig \(10\,000\;\text{kWh}\) oder \(3.60 \cdot 10^{10}\;\text{J}\)\\
\end{raster}
\fehler{\(P_\text{el} = \dfrac{P}{\eta} \approx 6.92\;\text{MW}\) (Wirkungsgrad umgekehrt): elektrisch 0; Energie folgerichtig \(\approx 13.8\;\text{MWh}\) zählt \(\to\) 3 von 4\,P.
Zeit in Sekunden mit Megawatt multipliziert und als «MWh» angegeben (\(36\,000\)): Energie 0 \(\to\) 3 von 4\,P.}
\end{lpaufgabe}

\begin{lpaufgabe}{G6}{Venus: \(S = 2601\;\text{W/m}^2\), Albedo \(0.77\); Erde}{5 P}
\begin{raster}
a) aufgenommen & 1\E & \((1 - a) \cdot \dfrac{S}{4} = 0.23 \cdot \dfrac{2601\;\text{W/m}^2}{4} \approx 150\;\text{W/m}^2\)\\
b) Temperatur & 1\E & \(\sigma \cdot T^4 = 149.6\;\text{W/m}^2\), \(T = \sqrt[4]{\dfrac{149.6\;\text{W/m}^2}{5.67 \cdot 10^{-8}\;\text{W/(m}^2\text{K}^4)}} \approx 227\;\text{K}\)\\
c) Unterschied & 1\B & Die dichte Atmosphäre (vor allem \(\text{CO}_2\)) lässt nur einen kleinen Teil der Wärmestrahlung der Oberfläche ins All: Die Bilanz ist erst bei viel höherer Oberflächentemperatur ausgeglichen (starker Treibhauseffekt)\\
d) Eis schmilzt & 1\B & betrifft die Aufnahme: \(1 - a\) wird grösser (\(0.71\) statt \(0.70\)), die Erde nimmt mehr auf, zuerst kommt mehr herein als hinaus — wärmer, bis die Abstrahlung wieder passt\\
d) mehr \(\text{CO}_2\) & 1\B & betrifft die Abstrahlung ins All: Weniger Wärmestrahlung gelangt hinaus, zuerst kommt mehr herein als hinaus — wärmer, bis die höhere Temperatur die Abstrahlung wieder ausgleicht\\
\end{raster}
\fehler{a) \(a \cdot \dfrac{S}{4} \approx 501\;\text{W/m}^2\) (zurückgeworfener Anteil): a) 0; b) folgerichtig \(\approx 307\;\text{K}\) zählt \(\to\) 4 von 5\,P.
a) Ohne Faktor 4 (\(\approx 598\;\text{W/m}^2\)): a) 0; b) folgerichtig \(\approx 320\;\text{K}\) \(\to\) 4 von 5\,P.
a) Albedo vergessen (\(\dfrac{S}{4} \approx 650\;\text{W/m}^2\)): a) 0; b) folgerichtig \(\approx 327\;\text{K}\) \(\to\) 4 von 5\,P.
b) Quadratwurzel statt vierter Wurzel: b) 0 \(\to\) 4 von 5\,P.
d) Richtung richtig, aber ohne Zuordnung zu Aufnahme oder Abstrahlung: diese Zeile 0. «\(\text{CO}_2\) erzeugt Wärme»: d) mehr \(\text{CO}_2\) 0 \(\to\) 4 von 5\,P.}
\end{lpaufgabe}

\needspace{14\baselineskip}
\section*{4 · Selbsteinschätzung}
\begin{tabularx}{\linewidth}{@{}l X@{}}\toprule
22 – 25 P & Die geprüften Teile sitzen. Wo du Punkte verloren hast: das Kapitel dieser Aufgabe nochmals (Zuordnung unten).\\
17 – 21 P & Den schwächsten Teil nochmals: Simulation und Übungen des Kapitels, in dem du die meisten Punkte verloren hast.\\
11 – 16 P & Zurück zu den Kapiteln aller Aufgaben, in denen du Punkte verloren hast.\\
0 – 10 P & Zurück zu Kapitel 1 und von dort der Reihe nach weiter.\\\midrule
\multicolumn{2}{@{}p{\linewidth}@{}}{\small\textbf{Aufgabe $\to$ Kapitel} (gilt für alle Bereiche): G1 $\to$ Kapitel 1, 5 (K1, K6) · G2 $\to$ Kapitel 1 (K2) · G3 $\to$ Kapitel 2, 3 (K3) · G4 $\to$ Kapitel 1, 4 (K2, K4) · G5 $\to$ Kapitel 5 (K6) · G6 $\to$ Kapitel 6 (K5)}\\\bottomrule
\end{tabularx}
\end{document}
